\documentclass{article}
%DIF LATEXDIFF DIFFERENCE FILE
%DIF DEL fixtures/synthetic/equation/001.tex                                 Thu Aug 27 03:19:23 2026
%DIF ADD test_output/equation_snapshot/.humanvoice/assembled/assembled.tex   Tue Sep 15 00:06:27 2026
\usepackage{amsmath}
%DIF PREAMBLE EXTENSION ADDED BY LATEXDIFF
%DIF UNDERLINE PREAMBLE %DIF PREAMBLE
\RequirePackage[normalem]{ulem} %DIF PREAMBLE
\RequirePackage{color}\definecolor{RED}{rgb}{1,0,0}\definecolor{BLUE}{rgb}{0,0,1} %DIF PREAMBLE
\providecommand{\DIFadd}[1]{{\protect\color{blue}\uwave{#1}}} %DIF PREAMBLE
\providecommand{\DIFdel}[1]{{\protect\color{red}\sout{#1}}} %DIF PREAMBLE
%DIF SAFE PREAMBLE %DIF PREAMBLE
\providecommand{\DIFaddbegin}{} %DIF PREAMBLE
\providecommand{\DIFaddend}{} %DIF PREAMBLE
\providecommand{\DIFdelbegin}{} %DIF PREAMBLE
\providecommand{\DIFdelend}{} %DIF PREAMBLE
\providecommand{\DIFmodbegin}{} %DIF PREAMBLE
\providecommand{\DIFmodend}{} %DIF PREAMBLE
%DIF FLOATSAFE PREAMBLE %DIF PREAMBLE
\providecommand{\DIFaddFL}[1]{\DIFadd{#1}} %DIF PREAMBLE
\providecommand{\DIFdelFL}[1]{\DIFdel{#1}} %DIF PREAMBLE
\providecommand{\DIFaddbeginFL}{} %DIF PREAMBLE
\providecommand{\DIFaddendFL}{} %DIF PREAMBLE
\providecommand{\DIFdelbeginFL}{} %DIF PREAMBLE
\providecommand{\DIFdelendFL}{} %DIF PREAMBLE
%DIF AMSMATHULEM PREAMBLE %DIF PREAMBLE
\makeatletter %DIF PREAMBLE
\let\sout@orig\sout %DIF PREAMBLE
\renewcommand{\sout}[1]{\ifmmode\text{\sout@orig{\ensuremath{#1}}}\else\sout@orig{#1}\fi} %DIF PREAMBLE
\makeatother %DIF PREAMBLE
%DIF COLORLISTINGS PREAMBLE %DIF PREAMBLE
\RequirePackage{listings} %DIF PREAMBLE
\RequirePackage{color} %DIF PREAMBLE
\lstdefinelanguage{DIFcode}{ %DIF PREAMBLE
%DIF DIFCODE_UNDERLINE %DIF PREAMBLE
  moredelim=[il][\color{red}\sout]{\%DIF\ <\ }, %DIF PREAMBLE
  moredelim=[il][\color{blue}\uwave]{\%DIF\ >\ } %DIF PREAMBLE
} %DIF PREAMBLE
\lstdefinestyle{DIFverbatimstyle}{ %DIF PREAMBLE
	language=DIFcode, %DIF PREAMBLE
	basicstyle=\ttfamily, %DIF PREAMBLE
	columns=fullflexible, %DIF PREAMBLE
	keepspaces=true %DIF PREAMBLE
} %DIF PREAMBLE
\lstnewenvironment{DIFverbatim}[1][]{\lstset{style=DIFverbatimstyle}}{} %DIF PREAMBLE
\lstnewenvironment{DIFverbatim*}[1][]{\lstset{style=DIFverbatimstyle,showspaces=true}}{} %DIF PREAMBLE
\lstset{extendedchars=\true,inputencoding=utf8}

%DIF END PREAMBLE EXTENSION ADDED BY LATEXDIFF

\begin{document}

This document tests equation extraction.

An inline equation: \DIFaddbegin \DIFadd{Einstein's mass-energy equivalence relation, }\DIFaddend $E = mc^2$\DIFaddbegin \DIFadd{, establishes that energy equals mass times the speed of light squared, revealing the fundamental convertibility between mass and energy.
}\DIFaddend 

\DIFdelbegin \DIFdel{A display equation}\DIFdelend \DIFaddbegin \DIFadd{The Pythagorean theorem states that in a right triangle, the sum of the squares of the two legs equals the square of the hypotenuse}\DIFaddend :
\begin{equation}
\label{eq:pythagorean}
a^2 + b^2 = c^2
\end{equation}
\DIFdelbegin %DIFDELCMD < 

%DIFDELCMD < %%%
\DIFdel{Another display equation without a label:
}\begin{displaymath}
\DIFdel{\int_{0}^{\infty} e^{-x^2} dx = \frac{\sqrt{\pi}}{2}
}\end{displaymath}%DIFAUXCMD
%DIFDELCMD < 

%DIFDELCMD < %%%
\DIFdel{The value }\DIFdelend \DIFaddbegin \DIFadd{where $a$ and $b$ represent the lengths of the two legs, and $c$ represents the length of the hypotenuse. This relationship, }\DIFaddend referenced by~\ref{eq:pythagorean}\DIFdelbegin \DIFdel{is fundamental}\DIFdelend \DIFaddbegin \DIFadd{, is fundamental to Euclidean geometry and underlies numerous results in mathematics and physics.
}

\DIFadd{The Gaussian integral evaluates to a closed form:
}\begin{equation*}
\DIFadd{\int_{0}^{\infty} e^{-x^2} dx = \frac{\sqrt{\pi}}{2}
}\end{equation*}
\DIFadd{demonstrating that the integral of the Gaussian function from zero to infinity equals $\sqrt{\pi}/2$}\DIFaddend .

\end{document}
